\section{Einleitung}
\label{sec:Einleitung}

% - Glossarquellen
% - Oberfragen und Unterfragen: Teilen auf jedenfall konstruktiver und nicht per Literaturrecherche
% - Konsistent halten der Hauptunterschriften
% - Die machen das so und so
% - Aufbau der Arbeit sollte Methodik und Aufbau heißen
% - 30-40 Seiten. Kein Fünftes mal erklären
% - Nicht über theoretische Grundlagen schreiben, welche man nicht braucht
% - Meetings über Anforderungen
% - Er hat eine Doktorarbeit über Performancemessung
% - Sicher und performante: Mandatenisolation -> schneller als die aktuelle Lösung
% - Schnell vs Durchsatz -> Starvation
% - Firecracker darf direkt rausfliegen.

\subsection{Forschungsfragen und Abgrenzung}
\label{sec:forschungsfragen}

Die Anbieter serverloser Plattformen haben ihre Ausführungsmodelle beschrieben und gegen die
Anforderungen des eigenen Betriebs begründet, etwa Amazon Web Services mit Firecracker
\autocite{agache_firecracker_2020} oder Cloudflare mit dem Isolationsmodell der
Workers-Plattform \autocite{cloudflare_security_2026}. Diese Arbeiten stellen jeweils eine
Lösung vor. Ein Ausführungsmodell, das für die mandantenisolierte Ausführung zur Laufzeit
erzeugter Funktionen konstruiert und anschließend gegen einen produktiven Ausgangszustand
gemessen wird, liegt dagegen nicht vor. Die folgende Oberfrage ist deshalb als
Gestaltungsfrage im Sinne der \ac{DSR} formuliert \autocite{osterle_memorandum_2010}. Sie
fragt nicht danach, welches der bekannten Modelle am besten abschneidet, sondern danach, wie
ein Modell beschaffen sein muss, das die genannten Eigenschaften zugleich erfüllt.

\begin{description}
  \item[FF] Wie muss ein Ausführungsmodell für nicht vertrauenswürdige, zur Laufzeit erzeugte
    JavaScript-Funktionen gestaltet sein, damit es in einer bestehenden Kubernetes-Plattform
    betreibbar bleibt, die Isolation zwischen Mandanten und die Zustandsfreiheit zwischen
    aufeinanderfolgenden Ausführungen sicherstellt und die Kaltstartlatenz gegenüber
    dem heutigen Betrieb der Fallstudie senkt?
\end{description}

Bezugspunkt der Frage ist der gemessene Ausgangszustand der Fallstudie und kein generischer
Vergleichswert. Die Oberfrage wird über fünf Unterfragen bearbeitet, die den vier Phasen des
Forschungsprozesses nach \textcite{osterle_memorandum_2010} folgen.

\begin{description}
  \item[U1] \label{rq:u1} Welche Anforderungen an ein solches Ausführungsmodell folgen aus dem
    Bedrohungsmodell der mandantenübergreifenden Ausführung und aus der Struktur des
    produktiven Skriptbestands der Fallstudie?

  \item[U2] \label{rq:u2} Welche Kombination aus Isolationsgrenze, Zuordnungsgranularität und
    Instanzverwaltung erfüllt diese Anforderungen, und wie lässt sie sich als lauffähiges
    Ausführungsmodell in die bestehende Kubernetes-Umgebung integrieren?

  \item[U3] \label{rq:u3} Hält das konstruierte Modell die geforderten Eigenschaften der
    Isolation und der Zustandsfreiheit ein, wenn es mit Prüffällen konfrontiert wird, die aus
    dem Bedrohungsmodell abgeleitet sind?

  \item[U4] \label{rq:u4} Welche Kaltstartlatenz, welche warme Latenz in den Perzentilen P50,
    P90, P95 und P99, welchen Durchsatz und welchen Speicherbedarf je Instanz erreicht der
    Prototyp unter beiden Lastprofilen im Vergleich zum gemessenen Ausgangszustand?

  \item[U5] \label{rq:u5} Welche über die Fallstudie hinaus übertragbaren Designprinzipien und
    welcher Migrationspfad folgen aus Entwurf und Evaluation?
\end{description}

U1 legt den Bewertungsmaßstab fest und stützt sich dabei nicht allein auf die Literatur. In
die Anforderungen geht eine Auswertung des produktiven Skriptbestands der Fallstudie ein, die
1.718 Mandanten mit 14.575 Skriptversionen danach klassifiziert, welchen Teil der
Node.js-Schnittstelle sie tatsächlich benötigen. U2 konstruiert das Artefakt. U3 und U4
bewerten es entlang der beiden Achsen des untersuchten Zielkonflikts, wobei U3 die
Sicherheitseigenschaften nicht aus den Herstellerangaben ableitet, sondern über ausgeführte
Prüffälle nachweist. U5 löst das Ergebnis von der Fallstudie und überführt es in
Designprinzipien.

Damit das Artefakt seine eigenen Kriterien verfehlen kann, werden die Zielwerte gemeinsam mit
dem Kriterienkatalog in Abschnitt~\ref{sec:kriterienkatalog} festgelegt, also vor dem Entwurf.
Ein dokumentiertes Verfehlen dieser Werte ist ein gültiges Ergebnis der Arbeit.

Nicht als Forschungsfrage geführt wird der Vergleich der bekannten Isolationstechnologien
untereinander. Ihre Isolations- und Startkosten sind in der Literatur beschrieben
\autocite{agache_firecracker_2020, young_true_2019, shillaker_faasm_2020}. Diese Angaben gehen
als Vorarbeit in die Vorauswahl in Abschnitt~\ref{sec:vorauswahl} ein und werden dort als
solche gekennzeichnet. Die Vorauswahl fällt dabei enger aus, als der Stand der Technik
zunächst nahelegt, weil der Ausgangszustand der Fallstudie mit gVisor bereits einen
User-Space-Kernel einsetzt, wie Abschnitt~\ref{sec:systemkontext} zeigt. Diese
Technologiefamilie gehört damit zur Referenz und nicht zu den zu prüfenden Alternativen.
Ebenfalls ausgeklammert bleibt eine erschöpfende Sicherheitsauditierung der einzelnen
Laufzeitumgebungen. Die Prüffälle aus U3 belegen das Verhalten des konstruierten Modells an
den Stellen, die das Bedrohungsmodell benennt, sie ersetzen keine systematische
Schwachstellensuche in den verwendeten Laufzeitumgebungen. Nicht betrachtet werden schließlich
die Qualität und die Korrektheit des erzeugten Codes sowie Maßnahmen, die das Modell selbst
gegen Beeinflussung über seine Eingaben härten \autocite{greshake_not_2023}. Gegenstand der
Arbeit ist allein die Frage, wie die Ausführung dieses Codes sicher und schnell erfolgen kann,
unabhängig davon, wie er zustande gekommen ist.